\chapter{A Thursday in Coleraine}

On Thursday 5 November 2026 I was at the sixth Northern Ireland High-Performance Computing User Conference, held in the Diamond Hall at Ulster University's Coleraine campus. It is a friendly, practical meeting. Researchers from Ulster and Queen's University Belfast who use the shared supercomputer, Kelvin2, come together to show what they have been running on it, swap tips about job queues and storage, and meet the people who keep the machines going. I was there to give a ninety-second lightning talk about my own PhD work: running the UK government's new Home Energy Model thousands of times on Kelvin2, to estimate what the next generation of Energy Performance Certificates might say about homes in London and Belfast.

The programme was a fair picture of where research computing sits in 2026. The first keynote came from the managing director of a company that builds data centres. The closing keynote, from an engineering manager at a data-storage manufacturer, was titled \emph{Powering the AI Era: The Future of Data and Advanced Manufacturing}. In between came talks on machine learning, drug design, vaccine candidates, solvents from seaweed, health inequalities, and the ethics of artificial intelligence. A great deal of the day, in other words, was about building more computing capacity, which is an industry with its own large and growing demand for electricity.

In the same block of lightning talks as mine was a talk titled \emph{Revisiting CO₂ emissions and global warming: Implications for society}. It presented a working paper, posted on SSRN~\citep{bhatta2026}, which analyses more than 70 million daily temperature readings from long-running weather stations and concludes that anthropogenic CO₂ is not the primary driver of global warming, and that current climate policy should be reassessed.

I want to be careful about how I respond. The author, Bibek Bhatta, is a researcher at Queen's University Belfast. He has done a lot of work with a large public dataset and has been open about his method, which is what makes it possible to check. This essay is not an attack on him. It is an attempt to do what scientists are supposed to do with any surprising claim: take the data seriously, rerun the analysis, see where the results differ, and explain why. Along the way I want to explain, as plainly as I can, how climate data work, because the mistakes in this case are ones that anyone can make if they come to the climate system from outside it.
